\chapter{Everyday use}
\label{ch:usage}

This chapter walks through the workflows you will run dozens of
times once the library is live: cataloging a new volume, searching
your collection, moving a volume, lending it out, taking it back,
and auditing what is on the shelves.

\section{The home page}\label{sec:home}

\begin{figure}[ht]
  \centering
  \includegraphics[width=0.95\textwidth]{01-home-desktop.png}
  \caption{The home page on a live install: search bar, genre
    filter chips, ``À traiter'' counters, ``Aperçu de la
    collection'' totals, and a recent-additions strip with cover
    thumbnails fetched through the metadata-provider chain.}
  \label{fig:home}
\end{figure}

The home page is the dashboard. From top to bottom:

\begin{itemize}
  \item \textbf{Scan field.}\index{scan field} A focused input that
        accepts an ISBN, a V-code or an L-code. The scanner-guard
        layer routes the input to the right page depending on its
        prefix.
  \item \textbf{Indicators row.} Recent additions, overdue loans,
        series with gaps, unshelved volumes — each click jumps to
        the filtered list.
  \item \textbf{Recent activity.} The last seven days of cataloging
        and returns.
\end{itemize}

\section{Cataloging a new title via barcode}\label{sec:catalog-scan}

\begin{enumerate}
  \item Focus the home-page scan field (it is selected by default
        on page load). Scan the ISBN / EAN-13 of the back cover.
  \item mybibli redirects to the catalog page, which displays a
        skeleton title with the barcode you scanned. A background
        task fires off requests against the metadata provider
        chain.
  \item Within a few seconds, the skeleton fills in: title,
        contributors, publisher, year, cover image. A small
        indicator confirms which provider resolved the title.
  \item Pick a \textbf{location}\index{location} from the dropdown
        (or scan the L-code of the destination shelf) and pick a
        \textbf{volume state}\index{volume state} (typically
        \emph{New}). Click \textbf{Add volume}.
  \item Print and stick a V-code label on the volume.
\end{enumerate}

If the metadata pipeline fails to resolve the ISBN, the title
appears with an \emph{Unresolved} badge. You can edit the title
fields manually and the volume is still created.

\begin{tip}
Cataloging in batches works best when the V-code sticker roll, the
scanner and the printer are all within arm's reach. Most of the time
is spent scanning, not waiting for mybibli — the metadata fetch is
asynchronous, so you do not have to wait for it to finish before
moving on to the next volume.
\end{tip}

\subsection*{Printing fresh label sheets}

Before printing a new sheet of V-code or L-code labels, check the
small line at the top of the catalog page (since 1.12.0, Librarian
and above): \emph{``Labels in use up to V0142~·~L0037 --- print new
label sheets starting after these numbers.''} The numbers include
items currently in the trash — a sticker that was ever printed and
stuck on a book is never proposed again, even if that volume was
deleted since.

\section{Searching the collection}\label{sec:search}

The \texttt{/search} page accepts:

\begin{itemize}
  \item full or partial \textbf{title};
  \item \textbf{contributor name} (matches on author by default;
        other contributor roles are scored lower);
  \item \textbf{ISBN / EAN-13} — exact match;
  \item \textbf{V-code} or \textbf{L-code} — exact match;
  \item \textbf{Dewey code} (for libraries that use Dewey).
\end{itemize}

Results are sorted by relevance. Click a result to open the
title page; click a contributor to filter the catalog by their
work; click a location label to see every volume currently in
that location.

\begin{tip}
You can ``search'' a single volume by scanning its V-code from any
page that has the home-page scan field — that includes the home
page, \texttt{/search} and \texttt{/loans}. Faster than typing.
\end{tip}

\subsection{Saved searches}\label{sec:saved-searches}

Recurring lookups (``comics without a cover'', ``uncategorized
titles'', ``everything overdue'') can be saved once and re-run with a
single click. On the home page, compose the search you want — type a
query, click a filter chip, choose a sort order — then open the
\textbf{Saved searches} dropdown next to the sort control and use
\emph{Save current search} to name and store it.

Saved searches are \textbf{shared across the whole instance}: every
librarian sees and can run the same set. To run one, open the dropdown
and click its name — the home page reloads with that query, filter and
sort applied. Each entry also offers \emph{Rename} and \emph{Delete}
(both confirm in a modal). Deleting is reversible: saving a new search
with the same name restores the old one.

\begin{note}
Saved searches capture the existing browse criteria (query, one filter
chip, sort column and direction) — they are not a free-form query
builder. The control is visible to librarians and administrators only.
\end{note}

\section{Management labels}\label{sec:management-labels}

A genre says what a book \emph{is}. A condition state says how a copy
has \emph{aged}. Neither says ``I am not sure this record is right,
come back to it'' --- and that sentence is the one a working catalogue
needs most. Since 1.18.0, \emph{management
labels}\index{label!management} carry it.

They are not the spine labels of section~\ref{sec:labels}: nothing is
printed, nothing is scanned, and an anonymous visitor never sees them.
They exist for the person maintaining the library.

\subsection*{Defining the vocabulary}

An admin creates the vocabulary once, in \emph{Administration
\textrightarrow{} Reference data \textrightarrow{} Labels}: ``To
check'', ``Re-read'', ``Damaged binding'', ``Lent informally'' ---
whatever your practice calls for. Labels are renamed and deleted from
the same place, and like every other reference value they are never
translated: what you type is what everyone sees, in any interface
language.

Keep the list short. A vocabulary of six labels is used; a vocabulary
of thirty is a second genre taxonomy nobody maintains.

\subsection*{Putting a label on something}

A librarian applies labels from the \emph{Labels} section of a title
page or of a single volume page, choosing from that one shared
vocabulary. The distinction matters:

\begin{itemize}
  \item on a \textbf{title}, the label is about the record --- ``the
        metadata here looks wrong'';
  \item on a \textbf{volume}, it is about that physical copy --- ``this
        one has a damaged binding'', which says nothing about the other
        two copies on the shelf.
\end{itemize}

\subsection*{Finding what carries a label}

\texttt{/labels} lists the vocabulary with a count against each entry.
Opening one shows what carries it, in two sections --- titles above
volumes --- each row linking through to its detail page. That is the
whole point: a label you cannot find your way back from is a note in a
drawer.

\begin{note}
Labels are internal. The handlers do not even run the queries for an
anonymous visitor, so nothing leaks through the public catalogue ---
including the fact that a label exists.
\end{note}

\section{Moving a volume to another location}\label{sec:move-volume}

\begin{enumerate}
  \item Open the volume page (search by V-code or click through from
        a title page).
  \item Click \textbf{Move}.
  \item Scan the L-code of the destination location, or pick it
        from the dropdown.
  \item Confirm.
\end{enumerate}

mybibli records the move in the volume's location history so you can
trace where a book has lived over time.

\section{Loans}\label{sec:loans}

\subsection*{Lending a volume}

\begin{enumerate}
  \item Open \texttt{/loans} and click \textbf{New loan}.
  \item Scan or pick the \textbf{borrower}\index{borrower}. If the
        borrower does not exist yet, click \emph{Create borrower}
        inline.
  \item Scan or pick the \textbf{volume} (V-code).
  \item Optionally set a custom \textbf{due date}; otherwise mybibli
        derives it from the overdue threshold (chapter
        \ref{ch:configuration}).
  \item Confirm.
\end{enumerate}

The volume's current location is recorded on the loan row so a
later return can restore it to the right shelf.

\subsection*{Returning a volume}

\begin{enumerate}
  \item Open \texttt{/loans} and find the active loan (filter on
        \emph{Active} or scan the V-code).
  \item Click \textbf{Return}.
  \item Pick the post-return state (\emph{Good}, \emph{Worn},
        \dots).
  \item Confirm. The volume snaps back to its pre-loan
        location automatically.
\end{enumerate}

\subsection*{Overdue loans}

Loans that exceed the overdue threshold turn red on the
\texttt{/loans} page and on the home-page \emph{Overdue} indicator.
You can run \texttt{/loans?filter=overdue} for a focused list.

\section{Shelf audit}\label{sec:audit}

An \emph{audit}\index{audit} is the routine ``walk the shelves and
confirm everything is where the database says it is''. Since 1.7.x
mybibli carries the workflow rather than leaving you to improvise it
around the search page.

The unit of work is a flag on a volume: \emph{to check}. You raise the
flag, walk the shelf, and lower it as each copy is confirmed.

\subsection*{Raising the flag}

\begin{itemize}
  \item \textbf{One volume} — on its detail page, \emph{Mark to check}.
        Useful when a reader mentions that something looks wrong, or
        when a scan surfaced a volume where it should not have been.
  \item \textbf{A whole shelf} — on a storage location's page,
        \emph{Mark all volumes for audit}. This is how a planned pass
        starts: flag the shelf on Sunday evening, walk it during the
        week.
\end{itemize}

A flagged volume carries a \emph{To check} chip wherever it is listed,
so the flag is visible without opening the audit view.

\subsection*{Walking the shelf}

The home dashboard shows a \emph{Volumes to check} indicator with a
count and a link; \texttt{/audit} is the list itself, sorted by
location then V-code — the order you physically walk in. Each row
carries a \emph{Checked} button: press it as you confirm the copy in
your hand, and the row leaves the list.

Two situations resolve themselves as you go:

\begin{itemize}
  \item A volume on the shelf that the list does not expect — scan it;
        mybibli says it belongs elsewhere and you move it
        (section~\ref{sec:move-volume}).
  \item A volume the list expects and the shelf does not hold — leave
        the flag raised. An unfinished audit list is a to-do list, not
        a failure; it is the only record of what you could not find.
\end{itemize}

\emph{Clear all flags} empties the list in one action. Use it when you
abandon a pass, not to finish one — a flag cleared without the volume
in your hand records something you did not verify.

\begin{tip}
A quarterly audit catches misplaced volumes early — the longer a
book has been on the wrong shelf, the harder it is to remember
where it should be.
\end{tip}

\section{Shelving workflow}\label{sec:shelving}

After a returns batch, the volumes pile up near the scanner
station. mybibli's \emph{shelving}\index{shelving} flow speeds up
putting them back:

\begin{enumerate}
  \item Pick up a volume and scan its V-code.
  \item mybibli shows the volume page with the original location
        prominently displayed.
  \item Walk the volume to that shelf. Scan the L-code on the shelf
        to confirm.
  \item Repeat.
\end{enumerate}

Confirming with the L-code is optional but recommended — it leaves
a fresh entry in the volume's location history that the next audit
can rely on.

\subsection{Undoing a scan action}\label{sec:scan-undo}

Barcode work is fast, and the occasional mis-scan happens — a volume
shelved at the wrong location, or a batch location activated by
mistake. When you shelve a volume or activate a batch location on the
catalog page, the feedback message shows an \textbf{Undo} button.
Click it within \textbf{30 seconds} and mybibli reverses that last
action:

\begin{itemize}
  \item Undoing a shelving puts the volume back at its previous
        location (or leaves it unshelved if it had no location before).
  \item Undoing a batch-location activation restores the previously
        active location.
\end{itemize}

Only the \emph{most recent} action can be undone, and the button
disappears once the 30-second window closes. Undo only reverses the
location change — it never deletes a title or volume you just
created. If the window has passed, move the volume the usual way
(\autoref{sec:move-volume}).

\section{Wish list}\label{sec:wishlist}

The \emph{wish list}\index{wish list} is a separate list of books
you'd like to buy, kept distinct from the owned catalog. It lives
under \texttt{/wishlist}\index{wishlist@/wishlist} (see the
navigation bar) and is available to Librarians and Admins.

\subsection{Adding an entry}

The \emph{Add to wish list} form has two modes:

\begin{itemize}
  \item \textbf{By ISBN.} Scan or type a 13-digit ISBN. mybibli
        looks the book up through the same metadata provider chain
        the catalog uses, then shows a preview card (title, author,
        publisher, year, cover). Click \emph{Add to wish list} to
        save. Cover images come from the resolved provider's URL
        directly (e.g.\ Open Library Covers or Google Books);
        downloading the cover locally is on the roadmap.
  \item \textbf{Free-form.} Use this when you don't have an ISBN
        handy (a self-published title, a recommendation jotted
        down at lunch). Only \emph{Title} is required; author,
        publisher and notes are optional.
\end{itemize}

\subsection{Browsing at the bookstore}

The list view collapses into a mobile-card stack on narrow
viewports — designed to scan quickly with one thumb at a
bookstore. Each card shows the title, the author (if known), the
ISBN (if any), the publisher, and the date you added it.

To print a paper copy or save a PDF for offline use, click
\emph{Print} at the top of the list. mybibli renders a clean
print-friendly page; use the browser's \emph{Print to PDF}
(File \string→ Print \string→ Save as PDF, or Ctrl/Cmd-P) to
generate a PDF you can take to the bookstore.

\subsection{Marking a book as bought}

When you buy a book, open its wish list entry (or tap \emph{Bought}
directly from the list row), confirm in the modal, and the entry
soft-deletes. The action is reversible from the Admin Trash panel
for 30 days.

\subsection{Auto-detect on catalog scan}

When you catalog a freshly bought book through the standard scan
flow (\texttt{/catalog}), mybibli checks whether its ISBN matches
an active wish list entry. If it does, an emerald banner appears
in the feedback area: \emph{This ISBN is on your wish list — Remove
it?} One click on the banner's \emph{Remove from wish list} button
opens the same \emph{Mark as bought} modal, and confirming removes
the entry. No silent deletions.

\section{What the library is worth}\label{sec:valuation}

Two optional amounts can be recorded against each volume: what it
\emph{cost} you (purchase price) and what it is \emph{worth} now
(current value), each with its own currency. Both are blank by default
and nothing forces you to fill them --- a library that is not an
investment can ignore this section entirely.

If you do record them, \texttt{/stats/value}\index{valuation} sums them
across the catalogue: a total per currency, then a breakdown by genre
and by series. Volumes with no value set are skipped rather than
counted as zero, so a partially-valued catalogue reports the truth
about the part you have valued rather than a diluted average.

\begin{itemize}
  \item The \textbf{default currency} for new valuations is set in
        \emph{Administration \textrightarrow{} System}; it is only a
        pre-fill, and any volume can carry another.
  \item The \textbf{home-dashboard card} showing the estimated total is
        \emph{off} by default. Turn it on in the same place if you want
        the figure in front of you every day --- some households do not.
\end{itemize}

\begin{note}
Insurance is the usual reason to keep these figures. If that is yours,
record the purchase price as well as the current value: the first is a
fact you can evidence, the second an estimate you will be asked to
justify.
\end{note}
